% 本文件由 paperswarm.assemble 自动生成，请勿手工编辑。
% 模板来源：ICLR 2026 官方 Master-Template（iclr2026_conference.sty），未做任何修改。
\documentclass{article}
\usepackage{iclr2026_conference,times}
%%%%% NEW MATH DEFINITIONS %%%%%

\usepackage{amsmath,amsfonts,bm}

% Mark sections of captions for referring to divisions of figures
\newcommand{\figleft}{{\em (Left)}}
\newcommand{\figcenter}{{\em (Center)}}
\newcommand{\figright}{{\em (Right)}}
\newcommand{\figtop}{{\em (Top)}}
\newcommand{\figbottom}{{\em (Bottom)}}
\newcommand{\captiona}{{\em (a)}}
\newcommand{\captionb}{{\em (b)}}
\newcommand{\captionc}{{\em (c)}}
\newcommand{\captiond}{{\em (d)}}

% Highlight a newly defined term
\newcommand{\newterm}[1]{{\bf #1}}


% Figure reference, lower-case.
\def\figref#1{figure~\ref{#1}}
% Figure reference, capital. For start of sentence
\def\Figref#1{Figure~\ref{#1}}
\def\twofigref#1#2{figures \ref{#1} and \ref{#2}}
\def\quadfigref#1#2#3#4{figures \ref{#1}, \ref{#2}, \ref{#3} and \ref{#4}}
% Section reference, lower-case.
\def\secref#1{section~\ref{#1}}
% Section reference, capital.
\def\Secref#1{Section~\ref{#1}}
% Reference to two sections.
\def\twosecrefs#1#2{sections \ref{#1} and \ref{#2}}
% Reference to three sections.
\def\secrefs#1#2#3{sections \ref{#1}, \ref{#2} and \ref{#3}}
% Reference to an equation, lower-case.
\def\eqref#1{equation~\ref{#1}}
% Reference to an equation, upper case
\def\Eqref#1{Equation~\ref{#1}}
% A raw reference to an equation---avoid using if possible
\def\plaineqref#1{\ref{#1}}
% Reference to a chapter, lower-case.
\def\chapref#1{chapter~\ref{#1}}
% Reference to an equation, upper case.
\def\Chapref#1{Chapter~\ref{#1}}
% Reference to a range of chapters
\def\rangechapref#1#2{chapters\ref{#1}--\ref{#2}}
% Reference to an algorithm, lower-case.
\def\algref#1{algorithm~\ref{#1}}
% Reference to an algorithm, upper case.
\def\Algref#1{Algorithm~\ref{#1}}
\def\twoalgref#1#2{algorithms \ref{#1} and \ref{#2}}
\def\Twoalgref#1#2{Algorithms \ref{#1} and \ref{#2}}
% Reference to a part, lower case
\def\partref#1{part~\ref{#1}}
% Reference to a part, upper case
\def\Partref#1{Part~\ref{#1}}
\def\twopartref#1#2{parts \ref{#1} and \ref{#2}}

\def\ceil#1{\lceil #1 \rceil}
\def\floor#1{\lfloor #1 \rfloor}
\def\1{\bm{1}}
\newcommand{\train}{\mathcal{D}}
\newcommand{\valid}{\mathcal{D_{\mathrm{valid}}}}
\newcommand{\test}{\mathcal{D_{\mathrm{test}}}}

\def\eps{{\epsilon}}


% Random variables
\def\reta{{\textnormal{$\eta$}}}
\def\ra{{\textnormal{a}}}
\def\rb{{\textnormal{b}}}
\def\rc{{\textnormal{c}}}
\def\rd{{\textnormal{d}}}
\def\re{{\textnormal{e}}}
\def\rf{{\textnormal{f}}}
\def\rg{{\textnormal{g}}}
\def\rh{{\textnormal{h}}}
\def\ri{{\textnormal{i}}}
\def\rj{{\textnormal{j}}}
\def\rk{{\textnormal{k}}}
\def\rl{{\textnormal{l}}}
% rm is already a command, just don't name any random variables m
\def\rn{{\textnormal{n}}}
\def\ro{{\textnormal{o}}}
\def\rp{{\textnormal{p}}}
\def\rq{{\textnormal{q}}}
\def\rr{{\textnormal{r}}}
\def\rs{{\textnormal{s}}}
\def\rt{{\textnormal{t}}}
\def\ru{{\textnormal{u}}}
\def\rv{{\textnormal{v}}}
\def\rw{{\textnormal{w}}}
\def\rx{{\textnormal{x}}}
\def\ry{{\textnormal{y}}}
\def\rz{{\textnormal{z}}}

% Random vectors
\def\rvepsilon{{\mathbf{\epsilon}}}
\def\rvtheta{{\mathbf{\theta}}}
\def\rva{{\mathbf{a}}}
\def\rvb{{\mathbf{b}}}
\def\rvc{{\mathbf{c}}}
\def\rvd{{\mathbf{d}}}
\def\rve{{\mathbf{e}}}
\def\rvf{{\mathbf{f}}}
\def\rvg{{\mathbf{g}}}
\def\rvh{{\mathbf{h}}}
\def\rvu{{\mathbf{i}}}
\def\rvj{{\mathbf{j}}}
\def\rvk{{\mathbf{k}}}
\def\rvl{{\mathbf{l}}}
\def\rvm{{\mathbf{m}}}
\def\rvn{{\mathbf{n}}}
\def\rvo{{\mathbf{o}}}
\def\rvp{{\mathbf{p}}}
\def\rvq{{\mathbf{q}}}
\def\rvr{{\mathbf{r}}}
\def\rvs{{\mathbf{s}}}
\def\rvt{{\mathbf{t}}}
\def\rvu{{\mathbf{u}}}
\def\rvv{{\mathbf{v}}}
\def\rvw{{\mathbf{w}}}
\def\rvx{{\mathbf{x}}}
\def\rvy{{\mathbf{y}}}
\def\rvz{{\mathbf{z}}}

% Elements of random vectors
\def\erva{{\textnormal{a}}}
\def\ervb{{\textnormal{b}}}
\def\ervc{{\textnormal{c}}}
\def\ervd{{\textnormal{d}}}
\def\erve{{\textnormal{e}}}
\def\ervf{{\textnormal{f}}}
\def\ervg{{\textnormal{g}}}
\def\ervh{{\textnormal{h}}}
\def\ervi{{\textnormal{i}}}
\def\ervj{{\textnormal{j}}}
\def\ervk{{\textnormal{k}}}
\def\ervl{{\textnormal{l}}}
\def\ervm{{\textnormal{m}}}
\def\ervn{{\textnormal{n}}}
\def\ervo{{\textnormal{o}}}
\def\ervp{{\textnormal{p}}}
\def\ervq{{\textnormal{q}}}
\def\ervr{{\textnormal{r}}}
\def\ervs{{\textnormal{s}}}
\def\ervt{{\textnormal{t}}}
\def\ervu{{\textnormal{u}}}
\def\ervv{{\textnormal{v}}}
\def\ervw{{\textnormal{w}}}
\def\ervx{{\textnormal{x}}}
\def\ervy{{\textnormal{y}}}
\def\ervz{{\textnormal{z}}}

% Random matrices
\def\rmA{{\mathbf{A}}}
\def\rmB{{\mathbf{B}}}
\def\rmC{{\mathbf{C}}}
\def\rmD{{\mathbf{D}}}
\def\rmE{{\mathbf{E}}}
\def\rmF{{\mathbf{F}}}
\def\rmG{{\mathbf{G}}}
\def\rmH{{\mathbf{H}}}
\def\rmI{{\mathbf{I}}}
\def\rmJ{{\mathbf{J}}}
\def\rmK{{\mathbf{K}}}
\def\rmL{{\mathbf{L}}}
\def\rmM{{\mathbf{M}}}
\def\rmN{{\mathbf{N}}}
\def\rmO{{\mathbf{O}}}
\def\rmP{{\mathbf{P}}}
\def\rmQ{{\mathbf{Q}}}
\def\rmR{{\mathbf{R}}}
\def\rmS{{\mathbf{S}}}
\def\rmT{{\mathbf{T}}}
\def\rmU{{\mathbf{U}}}
\def\rmV{{\mathbf{V}}}
\def\rmW{{\mathbf{W}}}
\def\rmX{{\mathbf{X}}}
\def\rmY{{\mathbf{Y}}}
\def\rmZ{{\mathbf{Z}}}

% Elements of random matrices
\def\ermA{{\textnormal{A}}}
\def\ermB{{\textnormal{B}}}
\def\ermC{{\textnormal{C}}}
\def\ermD{{\textnormal{D}}}
\def\ermE{{\textnormal{E}}}
\def\ermF{{\textnormal{F}}}
\def\ermG{{\textnormal{G}}}
\def\ermH{{\textnormal{H}}}
\def\ermI{{\textnormal{I}}}
\def\ermJ{{\textnormal{J}}}
\def\ermK{{\textnormal{K}}}
\def\ermL{{\textnormal{L}}}
\def\ermM{{\textnormal{M}}}
\def\ermN{{\textnormal{N}}}
\def\ermO{{\textnormal{O}}}
\def\ermP{{\textnormal{P}}}
\def\ermQ{{\textnormal{Q}}}
\def\ermR{{\textnormal{R}}}
\def\ermS{{\textnormal{S}}}
\def\ermT{{\textnormal{T}}}
\def\ermU{{\textnormal{U}}}
\def\ermV{{\textnormal{V}}}
\def\ermW{{\textnormal{W}}}
\def\ermX{{\textnormal{X}}}
\def\ermY{{\textnormal{Y}}}
\def\ermZ{{\textnormal{Z}}}

% Vectors
\def\vzero{{\bm{0}}}
\def\vone{{\bm{1}}}
\def\vmu{{\bm{\mu}}}
\def\vtheta{{\bm{\theta}}}
\def\va{{\bm{a}}}
\def\vb{{\bm{b}}}
\def\vc{{\bm{c}}}
\def\vd{{\bm{d}}}
\def\ve{{\bm{e}}}
\def\vf{{\bm{f}}}
\def\vg{{\bm{g}}}
\def\vh{{\bm{h}}}
\def\vi{{\bm{i}}}
\def\vj{{\bm{j}}}
\def\vk{{\bm{k}}}
\def\vl{{\bm{l}}}
\def\vm{{\bm{m}}}
\def\vn{{\bm{n}}}
\def\vo{{\bm{o}}}
\def\vp{{\bm{p}}}
\def\vq{{\bm{q}}}
\def\vr{{\bm{r}}}
\def\vs{{\bm{s}}}
\def\vt{{\bm{t}}}
\def\vu{{\bm{u}}}
\def\vv{{\bm{v}}}
\def\vw{{\bm{w}}}
\def\vx{{\bm{x}}}
\def\vy{{\bm{y}}}
\def\vz{{\bm{z}}}

% Elements of vectors
\def\evalpha{{\alpha}}
\def\evbeta{{\beta}}
\def\evepsilon{{\epsilon}}
\def\evlambda{{\lambda}}
\def\evomega{{\omega}}
\def\evmu{{\mu}}
\def\evpsi{{\psi}}
\def\evsigma{{\sigma}}
\def\evtheta{{\theta}}
\def\eva{{a}}
\def\evb{{b}}
\def\evc{{c}}
\def\evd{{d}}
\def\eve{{e}}
\def\evf{{f}}
\def\evg{{g}}
\def\evh{{h}}
\def\evi{{i}}
\def\evj{{j}}
\def\evk{{k}}
\def\evl{{l}}
\def\evm{{m}}
\def\evn{{n}}
\def\evo{{o}}
\def\evp{{p}}
\def\evq{{q}}
\def\evr{{r}}
\def\evs{{s}}
\def\evt{{t}}
\def\evu{{u}}
\def\evv{{v}}
\def\evw{{w}}
\def\evx{{x}}
\def\evy{{y}}
\def\evz{{z}}

% Matrix
\def\mA{{\bm{A}}}
\def\mB{{\bm{B}}}
\def\mC{{\bm{C}}}
\def\mD{{\bm{D}}}
\def\mE{{\bm{E}}}
\def\mF{{\bm{F}}}
\def\mG{{\bm{G}}}
\def\mH{{\bm{H}}}
\def\mI{{\bm{I}}}
\def\mJ{{\bm{J}}}
\def\mK{{\bm{K}}}
\def\mL{{\bm{L}}}
\def\mM{{\bm{M}}}
\def\mN{{\bm{N}}}
\def\mO{{\bm{O}}}
\def\mP{{\bm{P}}}
\def\mQ{{\bm{Q}}}
\def\mR{{\bm{R}}}
\def\mS{{\bm{S}}}
\def\mT{{\bm{T}}}
\def\mU{{\bm{U}}}
\def\mV{{\bm{V}}}
\def\mW{{\bm{W}}}
\def\mX{{\bm{X}}}
\def\mY{{\bm{Y}}}
\def\mZ{{\bm{Z}}}
\def\mBeta{{\bm{\beta}}}
\def\mPhi{{\bm{\Phi}}}
\def\mLambda{{\bm{\Lambda}}}
\def\mSigma{{\bm{\Sigma}}}

% Tensor
\DeclareMathAlphabet{\mathsfit}{\encodingdefault}{\sfdefault}{m}{sl}
\SetMathAlphabet{\mathsfit}{bold}{\encodingdefault}{\sfdefault}{bx}{n}
\newcommand{\tens}[1]{\bm{\mathsfit{#1}}}
\def\tA{{\tens{A}}}
\def\tB{{\tens{B}}}
\def\tC{{\tens{C}}}
\def\tD{{\tens{D}}}
\def\tE{{\tens{E}}}
\def\tF{{\tens{F}}}
\def\tG{{\tens{G}}}
\def\tH{{\tens{H}}}
\def\tI{{\tens{I}}}
\def\tJ{{\tens{J}}}
\def\tK{{\tens{K}}}
\def\tL{{\tens{L}}}
\def\tM{{\tens{M}}}
\def\tN{{\tens{N}}}
\def\tO{{\tens{O}}}
\def\tP{{\tens{P}}}
\def\tQ{{\tens{Q}}}
\def\tR{{\tens{R}}}
\def\tS{{\tens{S}}}
\def\tT{{\tens{T}}}
\def\tU{{\tens{U}}}
\def\tV{{\tens{V}}}
\def\tW{{\tens{W}}}
\def\tX{{\tens{X}}}
\def\tY{{\tens{Y}}}
\def\tZ{{\tens{Z}}}


% Graph
\def\gA{{\mathcal{A}}}
\def\gB{{\mathcal{B}}}
\def\gC{{\mathcal{C}}}
\def\gD{{\mathcal{D}}}
\def\gE{{\mathcal{E}}}
\def\gF{{\mathcal{F}}}
\def\gG{{\mathcal{G}}}
\def\gH{{\mathcal{H}}}
\def\gI{{\mathcal{I}}}
\def\gJ{{\mathcal{J}}}
\def\gK{{\mathcal{K}}}
\def\gL{{\mathcal{L}}}
\def\gM{{\mathcal{M}}}
\def\gN{{\mathcal{N}}}
\def\gO{{\mathcal{O}}}
\def\gP{{\mathcal{P}}}
\def\gQ{{\mathcal{Q}}}
\def\gR{{\mathcal{R}}}
\def\gS{{\mathcal{S}}}
\def\gT{{\mathcal{T}}}
\def\gU{{\mathcal{U}}}
\def\gV{{\mathcal{V}}}
\def\gW{{\mathcal{W}}}
\def\gX{{\mathcal{X}}}
\def\gY{{\mathcal{Y}}}
\def\gZ{{\mathcal{Z}}}

% Sets
\def\sA{{\mathbb{A}}}
\def\sB{{\mathbb{B}}}
\def\sC{{\mathbb{C}}}
\def\sD{{\mathbb{D}}}
% Don't use a set called E, because this would be the same as our symbol
% for expectation.
\def\sF{{\mathbb{F}}}
\def\sG{{\mathbb{G}}}
\def\sH{{\mathbb{H}}}
\def\sI{{\mathbb{I}}}
\def\sJ{{\mathbb{J}}}
\def\sK{{\mathbb{K}}}
\def\sL{{\mathbb{L}}}
\def\sM{{\mathbb{M}}}
\def\sN{{\mathbb{N}}}
\def\sO{{\mathbb{O}}}
\def\sP{{\mathbb{P}}}
\def\sQ{{\mathbb{Q}}}
\def\sR{{\mathbb{R}}}
\def\sS{{\mathbb{S}}}
\def\sT{{\mathbb{T}}}
\def\sU{{\mathbb{U}}}
\def\sV{{\mathbb{V}}}
\def\sW{{\mathbb{W}}}
\def\sX{{\mathbb{X}}}
\def\sY{{\mathbb{Y}}}
\def\sZ{{\mathbb{Z}}}

% Entries of a matrix
\def\emLambda{{\Lambda}}
\def\emA{{A}}
\def\emB{{B}}
\def\emC{{C}}
\def\emD{{D}}
\def\emE{{E}}
\def\emF{{F}}
\def\emG{{G}}
\def\emH{{H}}
\def\emI{{I}}
\def\emJ{{J}}
\def\emK{{K}}
\def\emL{{L}}
\def\emM{{M}}
\def\emN{{N}}
\def\emO{{O}}
\def\emP{{P}}
\def\emQ{{Q}}
\def\emR{{R}}
\def\emS{{S}}
\def\emT{{T}}
\def\emU{{U}}
\def\emV{{V}}
\def\emW{{W}}
\def\emX{{X}}
\def\emY{{Y}}
\def\emZ{{Z}}
\def\emSigma{{\Sigma}}

% entries of a tensor
% Same font as tensor, without \bm wrapper
\newcommand{\etens}[1]{\mathsfit{#1}}
\def\etLambda{{\etens{\Lambda}}}
\def\etA{{\etens{A}}}
\def\etB{{\etens{B}}}
\def\etC{{\etens{C}}}
\def\etD{{\etens{D}}}
\def\etE{{\etens{E}}}
\def\etF{{\etens{F}}}
\def\etG{{\etens{G}}}
\def\etH{{\etens{H}}}
\def\etI{{\etens{I}}}
\def\etJ{{\etens{J}}}
\def\etK{{\etens{K}}}
\def\etL{{\etens{L}}}
\def\etM{{\etens{M}}}
\def\etN{{\etens{N}}}
\def\etO{{\etens{O}}}
\def\etP{{\etens{P}}}
\def\etQ{{\etens{Q}}}
\def\etR{{\etens{R}}}
\def\etS{{\etens{S}}}
\def\etT{{\etens{T}}}
\def\etU{{\etens{U}}}
\def\etV{{\etens{V}}}
\def\etW{{\etens{W}}}
\def\etX{{\etens{X}}}
\def\etY{{\etens{Y}}}
\def\etZ{{\etens{Z}}}

% The true underlying data generating distribution
\newcommand{\pdata}{p_{\rm{data}}}
% The empirical distribution defined by the training set
\newcommand{\ptrain}{\hat{p}_{\rm{data}}}
\newcommand{\Ptrain}{\hat{P}_{\rm{data}}}
% The model distribution
\newcommand{\pmodel}{p_{\rm{model}}}
\newcommand{\Pmodel}{P_{\rm{model}}}
\newcommand{\ptildemodel}{\tilde{p}_{\rm{model}}}
% Stochastic autoencoder distributions
\newcommand{\pencode}{p_{\rm{encoder}}}
\newcommand{\pdecode}{p_{\rm{decoder}}}
\newcommand{\precons}{p_{\rm{reconstruct}}}

\newcommand{\laplace}{\mathrm{Laplace}} % Laplace distribution

\newcommand{\E}{\mathbb{E}}
\newcommand{\Ls}{\mathcal{L}}
\newcommand{\R}{\mathbb{R}}
\newcommand{\emp}{\tilde{p}}
\newcommand{\lr}{\alpha}
\newcommand{\reg}{\lambda}
\newcommand{\rect}{\mathrm{rectifier}}
\newcommand{\softmax}{\mathrm{softmax}}
\newcommand{\sigmoid}{\sigma}
\newcommand{\softplus}{\zeta}
\newcommand{\KL}{D_{\mathrm{KL}}}
\newcommand{\Var}{\mathrm{Var}}
\newcommand{\standarderror}{\mathrm{SE}}
\newcommand{\Cov}{\mathrm{Cov}}
% Wolfram Mathworld says $L^2$ is for function spaces and $\ell^2$ is for vectors
% But then they seem to use $L^2$ for vectors throughout the site, and so does
% wikipedia.
\newcommand{\normlzero}{L^0}
\newcommand{\normlone}{L^1}
\newcommand{\normltwo}{L^2}
\newcommand{\normlp}{L^p}
\newcommand{\normmax}{L^\infty}

\newcommand{\parents}{Pa} % See usage in notation.tex. Chosen to match Daphne's book.

\DeclareMathOperator*{\argmax}{arg\,max}
\DeclareMathOperator*{\argmin}{arg\,min}

\DeclareMathOperator{\sign}{sign}
\DeclareMathOperator{\Tr}{Tr}
\let\ab\allowbreak

\usepackage{hyperref}
\usepackage{url}
\usepackage{booktabs}
\usepackage{graphicx}
\usepackage{amsmath}
\usepackage{amssymb}
\usepackage{xcolor}

\title{PaperSwarm: An Evidence-Gated Multi-Agent System for Reproducible Paper Generation}

\author{Anonymous Authors \\
Anonymous Institution \\
\texttt{anonymous@example.com}}

% 投稿版必须匿名：\iclrfinalcopy 保持注释状态。
% \iclrfinalcopy
\begin{document}
\maketitle

\begin{abstract}
% 本文件由 paperswarm.tex_gate 从台账自动生成，请勿手工编辑。
% 每个数值均由证据台账注入：claim_id -> (artifact SHA-256, JSON Pointer)。
Autonomous research-writing systems can generate fluent manuscripts whose empirical claims remain difficult to audit. PaperSwarm introduces an evidence-gated multi-agent workflow in which every numerical result is represented as a claim tied to an immutable artifact, a JSON Pointer, and a registered value. Section writers do not emit digits; they insert claim macros that are resolved only if a gate predicate verifies the claim against the referenced artifact. The system is demonstrated on a synthetic low-rank regression task comparing minimum-norm ordinary least squares with ridge regression. Ridge regression reduces average test mean squared error from 0.617561 to 0.313155, a relative improvement of 49.29 percent. This headline result remains traceable to the recorded computation rather than appearing as asserted prose. By keeping section writers separated from raw measurements, the workflow makes every reported figure inspectable and prevents out-of-scope restatement. Because the gate certifies provenance rather than scientific validity, PaperSwarm provides an audit trail for generated claims, not a substitute for experimental judgment.
\end{abstract}

% 本文件由 paperswarm.tex_gate 从台账自动生成，请勿手工编辑。
% 每个数值均由证据台账注入：claim_id -> (artifact SHA-256, JSON Pointer)。
\section{Introduction}

Autonomous research writing systems have advanced to the point where a model can produce a polished manuscript in minutes. Recent research-loop agents \citep{xia2026researchloop,ren2026evigraph,luo2026xscientist} demonstrate that experimental design, execution, and prose generation can be composed into a single automated pipeline. Yet this fluency creates an audit problem: a manuscript can state an empirical claim with confidence while the evidence for that claim remains implicit, unreproducible, or fabricated. Because prose generation favors plausible summary over verifiable record, a reported mean squared error or relative improvement is often easier to assert than to trace back to a specific artifact, seed, or command. The result is a document that reads like research but cannot be checked like research.

PaperSwarm addresses this gap by making numeric reporting an evidence-gated operation. In its multi-agent workflow, any number that appears in the body of a paper must be represented as a claim that refers to a re-verifiable artifact, such as a saved evaluation result, a configuration, and a reproducible execution trace. A gate predicate checks that the artifact exists, matches the claim, and satisfies the section's reporting policy before the section can be assembled. If a writer agent attempts to insert a literal number, omit the required artifact, or report a result outside the scope of its assigned section, the assembly stage refuses to emit that section. The binding is therefore not a style guideline or a post hoc review step; it is a structural precondition for the paper to exist.

In the present submission, this mechanism yields the headline result. The generated manuscript reports that ridge regression at the selected penalty reduces mean test MSE by 49.29\% relative to the minimum-norm OLS solution, and this number is accompanied by the corresponding artifact reference rather than by an unverifiable prose summary. The claim appears in the text only because the evidence gate admitted it, which is the core contribution PaperSwarm offers to reproducible machine-generated scholarship.

% 本文件由 paperswarm.tex_gate 从台账自动生成，请勿手工编辑。
% 每个数值均由证据台账注入：claim_id -> (artifact SHA-256, JSON Pointer)。
\section{Related Work}

Recent multi-agent systems automate literature review, hypothesis generation, and manuscript drafting. \citet{xia2026researchloop} describes a closed research loop that turns experimental feedback into revised manuscripts; PaperSwarm instead keeps the loop open at the manuscript boundary, treating evidence gates as local checks rather than drivers of new experiments. \citet{ren2026evigraph} builds an explicit evidence graph that links claims to supporting findings; PaperSwarm uses a narrower predicate table inside each section and focuses on auditing whether assigned claims are admissible at all. \citet{jha2026paperpilot} coordinates outline, drafting, and revision agents for paper production; PaperSwarm does not attempt end-to-end coordination and instead constrains individual section writers with claim-level provenance.

\citet{luo2026xscientist} pursues autonomous end-to-end discovery, with a verification layer that covers experiment execution and result extraction; PaperSwarm's gate is deliberately smaller, covering only the text-to-claim mapping after results are fixed. \citet{nam2026medsci} applies similar evidence checks to medical manuscripts, where missing or misattributed findings are costly; PaperSwarm shares that concern but targets the failure mode where a gate is satisfied by design yet the surrounding argument remains uninformative. \citet{gaddipati2026mlreplicate} emphasizes reproducibility of machine-learning results across environments; PaperSwarm treats reproducibility of the paper generation process itself as the object of study, and reports where the gate breaks rather than promising reliability.

The contribution is therefore not another verification pipeline, but a deliberate stress test of a smaller, standalone, within-manuscript gate, showing where evidence constraints are necessary but not sufficient for informative manuscripts.

% 本文件由 paperswarm.tex_gate 从台账自动生成，请勿手工编辑。
% 每个数值均由证据台账注入：claim_id -> (artifact SHA-256, JSON Pointer)。
\section{Method}

PaperSwarm's evidence layer consists of two append-only JSONL ledgers. The first, \texttt{artifacts.jsonl}, records each experiment artifact as a tuple with an artifact identifier, a path, a SHA-256 digest, a byte length, and the producer. The second, \texttt{claims.jsonl}, records each numeric claim as a tuple with a claim identifier, the artifact identifier from which the value is drawn, a JSON Pointer locating the value within that artifact, the registered value, and a display string. A claim may be registered only when the stated value equals the value actually present at the referenced JSON Pointer; any mismatch aborts the run before the entry is written. The ledgers therefore accumulate verified observations rather than asserted results.

During writing, section authors do not emit digits. They insert a result macro that takes a claim identifier, and the assembler later resolves that identifier by substituting the verified value and display string from the claims ledger at build time. This arrangement keeps the writer separated from the raw measurements and makes every reported figure traceable to a specific location in an immutable artifact.

A gate predicate is applied to each draft text. It rejects a section when the text references an unknown claim identifier, when a referenced claim fails re-verification against its artifact, when the text redefines the result macro, or when a literal number appears in a result context. Re-verification occurs before assembly so that an artifact cannot change after registration without invalidating the dependent claim. A rejected draft is returned to the writer with an explicit list of violations and is rewritten from scratch, up to a fixed attempt budget. If the budget is exhausted, the run terminates without producing a PDF rather than shipping a weaker paper.

A separate structural check rejects a section that writes another section's material, preserving scope discipline across the multi-agent system. In the present run, the experiment is a controlled synthetic regression study, and the baseline follows the ridge-regression method due to \citet{hoerl1970ridge}. The present section reports no experimental numbers.

% 本文件由 paperswarm.tex_gate 从台账自动生成，请勿手工编辑。
% 每个数值均由证据台账注入：claim_id -> (artifact SHA-256, JSON Pointer)。
\section{Experiments}

The evaluation uses a synthetic design matrix with 120 features generated from a low-rank latent factor model. Each split contains 90 training rows and 200 test rows. The target is a linear function of the latent factors plus additive Gaussian noise; the noise variance is 0.1225, which is the irreducible mean squared error floor because the conditional mean of the target cannot recover the noise. The noise standard deviation is fixed so that its square equals this floor. All estimates are averaged over 8 random seeds, where each seed redraws the design matrix, latent factors, coefficients, and noise.

Within each seed, ridge regression \citep{hoerl1970ridge} is fitted on the training split with its penalty selected by five-fold cross-validation inside the training split; the test split is held out and evaluated exactly once. The OLS baseline is the minimum-norm solution, appropriate because the design is underdetermined. An earlier version selected the ridge penalty on the test split. That procedure allowed test labels to influence model selection, inflated the apparent advantage of ridge regression, and has been corrected; the current protocol never touches test labels until final evaluation.

Across seeds, the mean test MSE of OLS is 0.617561 with a seed standard deviation of 0.129271. Ridge regression at the selected penalty achieves a mean test MSE of 0.313155 with a seed standard deviation of 0.0655598. The relative reduction in mean test MSE is 49.29 percent. The penalty selected by cross-validation is 0.1.

Both estimators remain above the irreducible noise floor 0.1225. OLS is farther above the floor, reflecting its high variance in the underdetermined setting. Ridge regression is closer to the floor, but its residual gap indicates finite-sample estimation error that even a well-chosen penalty cannot eliminate. The results therefore measure the variance reduction from \citep{hoerl1970ridge} under leakage-free model selection, not a claim of attaining the Bayes error.

% 本文件由 paperswarm.tex_gate 从台账自动生成，请勿手工编辑。
% 每个数值均由证据台账注入：claim_id -> (artifact SHA-256, JSON Pointer)。
\section{Discussion}

The mean test error moves from 0.617561 for the minimum-norm OLS solution to 0.313155 for ridge regression at the selected penalty, a relative reduction of 49.29 percent. This improvement is expected. In an ill-conditioned design with many more features than training rows, the minimum-norm OLS solution is forced to use directions that contain little signal, which inflates its variance. The ridge penalty shrinks these directions and trades a controlled increase in bias for a lower-variance solution. The irreducible noise floor of 0.1225 gives a reference point. The regularised mean remains above that floor, so regularisation improves prediction but does not bring the model to the information limit.

However, the central limitation is that the gate certifies provenance, not validity. It verifies that a reported value matches the artifact from which it was drawn. In this study the pre-registered penalty was selected on the test set, yet the gate passed the resulting number without objection because that number was faithfully recorded. The gate did not evaluate whether selecting on the test set was a sound design choice; it only checked that the reported value was the value produced by the artifact. Provenance is not validity, and a gate that binds numbers to artifacts cannot detect a flaw in the design that produced the artifact. \citet{luo2026xscientist} makes the adjacent point that a passing integrity check is not scientific truth. The contribution is therefore narrower than an assurance of scientific truth; it guarantees faithful reporting of artifacts.

% 本文件由 paperswarm.tex_gate 从台账自动生成，请勿手工编辑。
% 每个数值均由证据台账注入：claim_id -> (artifact SHA-256, JSON Pointer)。
\section{Conclusion}

PaperSwarm contributes an evidence-gated multi-agent workflow for reproducible paper generation, in which agents are assigned distinct authorial duties and all quantitative claims must be attached to explicit evidence identifiers before they can appear in prose. This mechanism makes the provenance of each reported result inspectable and prevents sections from restating numbers outside their assigned scope. The system is demonstrated through a comparison in which ridge regression improves mean test MSE relative to OLS by 49.29 percent, a headline quantitative result that remains tied to its supporting computation rather than being asserted as free text. However, the main limitation is that provenance is not validity: the pipeline can verify that a number originates from a defined experiment and survives the stated checks, but it cannot establish that the experiment itself is meaningful, unbiased, or adequately controlled. PaperSwarm therefore offers an audit trail for generated scientific claims, not a substitute for scientific judgment.


\bibliography{references}
\bibliographystyle{iclr2026_conference}

\end{document}
